% Restyle of fig-roles to the tikz-diagram-craft v2 standard (v1 names and
% raw colours migrated to roles; three lifelines, four messages, unchanged
% layout, ch4.md 4.14).
%
% NOTE for the author: the atlas row (I/fig:roles) prescribes a
% "differentiated institutional route" encoding operator/principal,
% multitude, enforcement, legibility and mutual covenant explicitly; this
% fragment draws a simpler three-lifeline grant-and-gate sequence (operator,
% daemon, agents) that matches the chapter's own roles paragraph and
% caption. The gap between the atlas's fuller cast and this figure's three
% lifelines is flagged for the author, not resolved here.
%
% Reader question: what grants authority, who exercises it, who is
% governed, and how is override returned?
% Claim: one grant lets many agents act, but the daemon -- not the agents --
% gates every act, and only the operator can revoke.
% Evidence: the grant-act-gate-override sequence, whitepaper/
% legible-swarm.tex, the roles section [internal].
% Must distinguish: the one-time grant (message 1) from the repeated
% gate (messages 2-3, per act); the override (message 4) as the one edge
% that bypasses the agents entirely.
% Grammar chosen: sequence with lifelines and numbered messages, no
% crossing -- the claim is entirely about who sends what to whom, in what
% order, which this grammar reads directly.
% Grammar rejected: a state machine would suggest the daemon itself changes
% state per act; the daemon's role here is a standing check, not a state.
% Five-second test: which edge lets the operator bypass the agents?
\begin{figure}[H]
\centering
\pdfigurehue{pdteal}%
\begin{tikzpicture}[pd figure,x=0.82cm,y=1cm]
  \node[pd state] (op) at (0,8.0) {\textbf{operator}\\[1pt]\footnotesize author};
  \node[pd state] (dm) at (5.0,8.0) {\textbf{daemon}\\[1pt]\footnotesize actor};
  \node[pd state] (ag) at (10.0,8.0) {\textbf{agents}\\[1pt]\footnotesize multitude};

  \draw[pd guide] (0,7.35) -- (0,0.3);
  \draw[pd guide] (5.0,7.35) -- (5.0,0.3);
  \draw[pd guide] (10.0,7.35) -- (10.0,0.3);

  % 1. operator grants the daemon a scoped grant g
  \draw[pd focus arrow] (0,6.1) -- (5.0,6.1);
  \node[pd focus label,anchor=south] at (2.5,6.2) {scoped grant g};

  % 2. an agent's act arrives at the daemon under g
  \draw[pd arrow] (10.0,4.6) -- (5.0,4.6);
  \node[pd label,anchor=south] at (7.5,4.7) {act a1 under g};

  % 3. the daemon gates or escalates that act
  \draw[pd arrow] (5.0,3.1) -- (10.0,3.1);
  \node[pd label,anchor=south] at (7.5,3.2) {gate / escalate};

  % 4. the override: the one edge from operator to daemon
  \draw[pd breach arrow] (0,1.6) -- (5.0,1.6);
  \node[pd label,text=pderror,anchor=south] at (2.5,1.7) {override: revoke g};

  \node[pd note,anchor=north,align=center,text width=7.6cm] at (5.0,-0.1)
    {\hyphenpenalty=10000\exhyphenpenalty=10000
     one grant; the daemon gates every act;\\only the operator overrides};
\end{tikzpicture}
\caption{The operator grants once; the daemon gates or escalates each subsequent act. Operator override bypasses agents. A live fleet repeats the act-and-check pair, not the grant. [internal]}
\label{fig:roles}
\end{figure}
